
\documentclass[11pt]{article}
\usepackage{series}

\renewcommand\SeriesShort{Delta Projection Execution Program}
\renewcommand\PartRoman{I}
\renewcommand\PartArabic{1}
\renewcommand\PartsInSeries{I}

\title{Admissibility-Time Bounds from Fixed--Point Damping\\
\large Exact Exit Times, Certified Proper-Time Filters, and the MTT Source Boundary}
\author{Peter Nero}
\date{July 2026, Version 1}

\begin{document}
\maketitle

\begin{abstract}
The preceding delta-projection papers established a structural dictionary: Dirac delta functions arise as singular limits of finite kernels, filters, shells, or selection operations. A natural objection is that this does not yet select a physical finite kernel. This paper isolates what fixed-point damping does and does not determine. If the incoherent sector of an evolve--project system obeys the semigroup estimate
\[
\|Q\Phi_t Q\| \le C_Q e^{-\lambda_\ast t},
\]
and admissibility requires residual leakage below \(\varepsilon_{\rm adm}\), then
\[
\tau_{\rm cert}=\frac{1}{\lambda_\ast}\log\frac{C_Q}{\varepsilon_{\rm adm}}
\]
is the earliest time certified by that bound. It is not, in general, the exact first admissible time. We define the exact exit time directly and prove equality in the reducing self-adjoint case, where
\(\|Qe^{-tA}Q\|=e^{-t\inf\sigma(A|_{\operatorname{Ran}Q})}\).
For supplied operator, projector, acceptance window, and tolerance data, the corresponding filtered operator is
\[
K_{\rm adm}(x,y)
=
\big\langle x\big|P\,\chi(A)\,e^{-\tau A}\,\chi(A)\,P\big|y\big\rangle,
\]
where \(A\) is the linearized stabilization operator. We derive an exact error decomposition and a circle benchmark. The result is a conditional execution bridge, not a source theorem: MTT still has to select \(A\), \(P\), \(\chi\), and \(\varepsilon_{\rm adm}\) in the physical sector before this filter can replace a regulator or predict a coherence scale.
\end{abstract}

\section*{Version 1 Revision Note}
\begin{description}
\item[Supersedes] The unversioned April 2026 manuscript on fixed-point damping and
proper-time kernels.
\item[Reason] The earlier theorem called a sufficient semigroup threshold the minimal
physical admissibility time and described the resulting kernel as MTT-selected.
\item[Resolution] This version distinguishes the exact exit time from the earliest time
certified by an exponential upper bound, proves equality in the reducing self-adjoint case,
and makes the operator, projector, spectral window, tolerance, and physical intertwiner
explicit source obligations.
\item[Retained result] The semigroup threshold formula, functional-calculus filter, error
decomposition, and circle benchmark remain valid at their corrected conditional tier.
\item[Remaining boundary] MTT geometry must select the physical source data before the
filter predicts a coherence scale or modifies a propagator.
\end{description}

\section{Purpose and Claim Discipline}

Earlier papers in the delta-projection program established the forward singular-limit dictionary
\[
K_\epsilon \longrightarrow \delta .
\]
That result is mathematically legitimate but, by itself, does not constitute new physical prediction. The decisive reverse problem is:
\[
\text{MTT fixed-point data}
\Longrightarrow
(A,P,\chi,\tau)
\Longrightarrow
K_{\rm coh}
\Longrightarrow
\text{finite correction}.
\]

This paper solves a narrower part of that reverse problem. It converts quantified damping and a supplied admissibility tolerance into either an exact exit time or a certified upper bound. Functional calculus then determines a filter from the supplied operator data. The result does not select those inputs from MTT geometry.

\subsection*{What is proved}

We prove the following statements in a Hilbert-space fixed-point setting:

\begin{enumerate}
\item A residual incoherent leakage bound determines the earliest time certified by that bound.
\item The exact first admissible time is characterized separately and is computed in the reducing self-adjoint case.
\item Either time parameter determines a proper-time filter once the choice between exact and conservative execution is declared.
\item Together with a declared coherent projector \(P\) and spectral window \(\chi(A)\), this fixes the conditional kernel \(K_{\rm adm}\).
\item Error terms split into damping leakage, spectral-window truncation, and proper-time smoothing.
\end{enumerate}

\subsection*{What is not yet proved}

This paper does not derive the Standard Model, the numerical physical value of \(\varepsilon_{\rm adm}\), a unique universal \(\tau\), or the selected physical operator/window. The result is regime-relative and conditional on those data.

\section{Fixed-Point Data}

Let \(\mathcal H\) be a Hilbert space of modal perturbations around a coherent fixed point \(\Psi_\ast\). Let
\[
\mathcal H=\mathcal H_{\rm coh}\oplus \mathcal H_{\rm inc}
\]
with orthogonal projectors
\[
P:\mathcal H\to\mathcal H_{\rm coh},
\qquad
Q=I-P.
\]

Let \(A\ge0\) be a self-adjoint nonnegative operator representing the linearized stabilization operator around the fixed point. In the simplest gradient-flow case, \(A\) is the Hessian or linearized dissipative generator of the closure/admissibility functional. We assume the incoherent sector has a positive gap:
\[
A|_{\operatorname{Ran}Q}\ge \lambda_\ast I,
\qquad \lambda_\ast>0.
\]

We also assume the coherent projection is compatible with the linearized operator:
\[
[P,A]=0
\]
on the relevant domain. This assumption is automatic when \(P\) is a spectral/Riesz projector of \(A\). If \(P\) is only an admissible sector projector not exactly spectral for \(A\), the results below apply after replacing commutation by the corresponding bounded commutator estimates; that perturbative extension is not needed for the core theorem.

Let \(\Phi_t\) denote the linearized or locally linearized damping flow. In the linear model,
\[
\Phi_t=e^{-tA}.
\]
In the nonlinear local fixed-point model, the same estimates hold on an admissible slab after replacing \(C_Q\) and \(\lambda_\ast\) by the constants in the local semigroup/contractivity estimate.

\section{Admissibility Tolerance and Minimal Projection Time}

The key new input is not a regulator. It is an admissibility threshold.
Throughout this section, \(t\) and \(\tau\) denote heat/proper-time parameters for the semigroup \(e^{-tA}\), not ordinary Lorentzian clock time. If \(A\) is Laplace-type, then its eigenvalues have dimensions of energy squared in units \(\hbar=c=1\), and \(t\) or \(\tau\) has dimensions of inverse energy squared.

\begin{definition}[Admissible residual leakage]
Let \(\varepsilon_{\rm adm}\in(0,1)\) be a fixed tolerance for incoherent leakage. We say that the projected stabilization step of duration \(t\) is \(\varepsilon_{\rm adm}\)-admissible if
\[
\|Q\Phi_t Q\|_{\mathcal H\to\mathcal H}
\le \varepsilon_{\rm adm}.
\]
\end{definition}

The tolerance \(\varepsilon_{\rm adm}\) encodes the maximum residual noncoherent contribution allowed inside the effective description. It is not a universal constant in this paper; it is regime data. In an execution-level sector paper, it must be derived from closure-strain, detector resolution, observed error tolerance, or basin-separation data.

\begin{assumption}[Incoherent damping bound]
There exist constants \(C_Q\ge1\) and \(\lambda_\ast>0\) such that
\[
\|Q\Phi_tQ\|\le C_Qe^{-\lambda_\ast t}
\qquad (t\ge0).
\]
\end{assumption}

\begin{definition}[Exact admissibility exit time]
For \(0<\varepsilon_{\rm adm}<1\), define
\[
\tau_{\rm exit}
:=
\inf\{t\ge0:\|Q\Phi_tQ\|\le\varepsilon_{\rm adm}\}.
\]
When the leakage norm is nonincreasing, every \(t\ge\tau_{\rm exit}\) is admissible.
\end{definition}

\begin{theorem}[Certified admissibility time]
Under the incoherent damping bound, any time \(t\) satisfying
\[
t\ge \frac{1}{\lambda_\ast}\log\frac{C_Q}{\varepsilon_{\rm adm}}
\]
is \(\varepsilon_{\rm adm}\)-admissible. The earliest time certified by this bound is
\[
\boxed{
\tau_{\rm cert}
=
\frac{1}{\lambda_\ast}\log\frac{C_Q}{\varepsilon_{\rm adm}} .
}
\]
If the leakage norm is nonincreasing, then
\[
\tau_{\rm exit}\le\tau_{\rm cert}.
\]
\end{theorem}

\begin{proof}
The damping estimate gives
\[
\|Q\Phi_tQ\|\le C_Qe^{-\lambda_\ast t}.
\]
Requiring the right-hand side to be at most \(\varepsilon_{\rm adm}\) gives
\[
C_Qe^{-\lambda_\ast t}\le \varepsilon_{\rm adm}.
\]
Taking logarithms yields
\[
-\lambda_\ast t\le \log \varepsilon_{\rm adm}-\log C_Q
\]
and hence
\[
t\ge \lambda_\ast^{-1}\log(C_Q/\varepsilon_{\rm adm}).
\]
This is the earliest time at which the stated upper bound itself falls below the tolerance.
The actual norm may cross earlier, so the result gives
\(\tau_{\rm exit}\le\tau_{\rm cert}\), not equality in general. \(\square\)
\end{proof}

\begin{theorem}[Exact exit time for a reducing self-adjoint semigroup]
Assume \(\Phi_t=e^{-tA}\), \(A\ge0\) is self-adjoint, \(Q\) reduces \(A\), and
\[
\lambda_Q:=\inf\sigma(A|_{\operatorname{Ran}Q})>0.
\]
Then
\[
\|Qe^{-tA}Q\|=e^{-t\lambda_Q}
\]
and hence
\[
\boxed{\tau_{\rm exit}=\lambda_Q^{-1}\log(1/\varepsilon_{\rm adm}).}
\]
\end{theorem}

\begin{proof}
Because \(Q\) reduces \(A\), the restriction \(A_Q=A|_{\operatorname{Ran}Q}\) is
self-adjoint. The spectral theorem gives
\[
\|e^{-tA_Q}\|=\sup_{\lambda\in\sigma(A_Q)}e^{-t\lambda}
=e^{-t\inf\sigma(A_Q)}.
\]
Solving \(e^{-t\lambda_Q}\le\varepsilon_{\rm adm}\) gives the formula. \(\square\)
\end{proof}

\begin{remark}[Regulator choice versus conditional selection]
A heat-kernel regulator normally chooses \(\tau\) as a computational cutoff. Here a supplied
tolerance and a selected damping operator determine an exact or certified time before a loop
integral is considered. This removes freedom only after those inputs have been sourced.
\end{remark}


\paragraph{What has and has not been fixed.}
This result fixes an exact or conservative proper-time scale once the damping data and admissibility tolerance are specified:
\[
(A,Q,\varepsilon_{\rm adm})\quad\Longrightarrow\quad \tau_{\rm exit},
\qquad
(C_Q,\lambda_\ast,\varepsilon_{\rm adm})\quad\Longrightarrow\quad \tau_{\rm cert}.
\]
Thus \(\tau\) is no longer an arbitrary Gaussian width conditional on the inputs. The remaining execution-level task is sharper: in each physical sector one must derive or constrain \(\varepsilon_{\rm adm}\), together with \(A\), \(P\), \(\chi\), \(C_Q\), and \(\lambda_\ast\), from selected source data. Detector resolution or measured error tolerance may define an effective filter, but cannot establish a first-principles MTT scale.

\section{Canonical Admissible Kernel}

Let \(\chi:[0,\infty)\to[0,1]\) be a bounded Borel spectral acceptance window. We assume \(\chi(A)\) commutes with \(P\), which follows if \([P,A]=0\). The window \(\chi\) encodes retained coherent bandwidth. Typical choices are:

\[
\chi(\lambda)=\mathbf 1_{\lambda\le\Lambda_{\rm coh}^2}
\]
for a sharp spectral window, or a smooth cutoff satisfying
\[
\chi(\lambda)\approx1 \quad (\lambda\le\Lambda_{\rm coh}^2),
\qquad
\chi(\lambda)\approx0 \quad (\lambda\ge\Lambda_{\rm inc}^2).
\]

\begin{definition}[Canonical admissible MTT filter]
The admissible fixed-point filter is
\[
B_{\tau}
=
P\,\chi(A)\,e^{-\tau A}\,\chi(A)\,P .
\]
When the corresponding operator has an integral kernel, we define
\[
\boxed{
K_{\tau}(x,y)
=
\langle x|B_{\tau}|y\rangle .
}
\]
\end{definition}

This is the central execution object. It replaces the arbitrary \(K_\epsilon\) of the dictionary papers.

\begin{theorem}[Conditional uniqueness of the filtered operator]
For declared regime data
\[
(A,P,\chi,\tau),
\]
the filtered operator \(B_\tau\) is uniquely determined by functional calculus. If \(\tau\)
is chosen as \(\tau_{\rm exit}\) or \(\tau_{\rm cert}\), it is fixed by the corresponding declared
damping rule.
\end{theorem}

\begin{proof}
Given \(A\), the spectral theorem uniquely defines \(\chi(A)\) and \(e^{-\tau A}\). Given \(P\), the product
\[
B_\tau=P\chi(A)e^{-\tau A}\chi(A)P
\]
is therefore uniquely determined as a bounded operator. If an integral kernel representation exists in the chosen function space, that kernel is the Schwartz kernel of \(B_\tau\), hence is uniquely determined as a distribution and, under smoothing hypotheses, as a function. \(\square\)
\end{proof}

\section{Error Decomposition}

The full identity \(I\) is not generally the correct target. In MTT, \(P\) is a sector identity: it is the identity on the retained coherent sector, not on the full upstream space.

Let
\[
B_\tau=P\chi(A)e^{-\tau A}\chi(A)P.
\]
For \(f\in\mathcal H\),
\[
f-B_\tau f
=
(I-P)f
+
P(I-\chi(A)^2)Pf
+
P\chi(A)\big(I-e^{-\tau A}\big)\chi(A)Pf .
\]
This is the key execution decomposition.

\begin{theorem}[Three-source correction bound]
Assume \([P,A]=0\) and \(0\le\chi\le1\). Then
\[
\|f-B_\tau f\|
\le
\|(I-P)f\|
+
\|(I-\chi(A)^2)Pf\|
+
\|\chi(A)(I-e^{-\tau A})\chi(A)Pf\|.
\]
If \(f\in\operatorname{Ran}P\) and \(\chi(A)f=f\), then
\[
\|f-B_\tau f\|
\le
\|(I-e^{-\tau A})f\|.
\]
If additionally the retained spectral support of \(f\) lies in \([0,\Lambda^2]\), then
\[
\|f-B_\tau f\|
\le
\left(1-e^{-\tau\Lambda^2}\right)\|f\|.
\]
For \(\tau\Lambda^2\ll1\),
\[
\|f-B_\tau f\|
\le
\tau\Lambda^2\|f\|+O(\tau^2\Lambda^4)\|f\|.
\]
\end{theorem}

\begin{proof}
The algebraic decomposition follows by adding and subtracting \(Pf\), \(P\chi(A)^2Pf\), and \(P\chi(A)e^{-\tau A}\chi(A)Pf\), using \(P^2=P\) and commutation. The norm bound follows by the triangle inequality. If \(f\in\operatorname{Ran}P\) and \(\chi(A)f=f\), the first two terms vanish. On spectral support \([0,\Lambda^2]\), functional calculus gives
\[
\|(I-e^{-\tau A})f\|
\le
\sup_{0\le\lambda\le\Lambda^2}|1-e^{-\tau\lambda}|\|f\|
=
(1-e^{-\tau\Lambda^2})\|f\|.
\]
The small-\(\tau\) expansion follows from \(1-e^{-x}=x+O(x^2)\). \(\square\)
\end{proof}

\begin{center}
\begin{tabular}{lll}
\hline
Term & Formula & Meaning\\
\hline
Projection loss & \((I-P)f\) & discarded noncoherent sector\\
Window loss & \((I-\chi(A)^2)Pf\) & modes outside acceptance band\\
Proper-time smoothing & \((I-e^{-\tau A})\chi(A)Pf\) & finite coherence scale\\
\hline
\end{tabular}
\end{center}

\section{Effective Scale}

In flat or locally flat sectors where \(A\sim -\Delta\), spectral values scale like \(\lambda\sim k^2\). The damping factor becomes
\[
e^{-\tau k^2}.
\]
The characteristic scale at which damping becomes order one is
\[
\tau k^2\sim1.
\]
Thus
\[
\boxed{
\Lambda_{\rm eff}\sim \tau^{-1/2}.
}
\]
Using the certified damping expression \(\tau=\tau_{\rm cert}\),
\[
\boxed{
\Lambda_{\rm eff}
\sim
\left(
\frac{\lambda_\ast}{\log(C_Q/\varepsilon_{\rm adm})}
\right)^{1/2}.
}
\]
This is an explicit route from declared fixed-point damping data to a finite-width filter
scale. It becomes a physical correction scale only after source selection and matching.

\section{Worked Model: Circle Sector}

Let
\[
X=S^1_R
\]
be the circle of radius \(R\), with coordinate \(\theta\in[0,2\pi)\). Take
\[
A=-\frac{1}{R^2}\partial_\theta^2.
\]
The eigenfunctions are
\[
\phi_n(\theta)=\frac{1}{\sqrt{2\pi R}}e^{in\theta},
\qquad
\lambda_n=\frac{n^2}{R^2},
\qquad n\in\mathbb Z.
\]

Assume the coherent regime retains a symmetric band \(|n|\le N_{\rm coh}\). Then
\[
P=P_{N_{\rm coh}}
=
\sum_{|n|\le N_{\rm coh}}|\phi_n\rangle\langle\phi_n|.
\]
For simplicity take \(\chi=1\) on the retained band. Let the first discarded mode have
\[
\lambda_\ast=\frac{(N_{\rm coh}+1)^2}{R^2}.
\]
Assume \(C_Q=1\). Then
\[
\tau_{\rm adm}
=
\frac{R^2}{(N_{\rm coh}+1)^2}
\log\frac{1}{\varepsilon_{\rm adm}}.
\]

The canonical kernel is
\[
K_{\rm adm}(\theta,\theta')
=
\frac{1}{2\pi R}
\sum_{|n|\le N_{\rm coh}}
\exp\left[
-\frac{\tau_{\rm adm}n^2}{R^2}
\right]
e^{in(\theta-\theta')}.
\]
Substituting \(\tau_{\rm adm}\),
\[
K_{\rm adm}(\theta,\theta')
=
\frac{1}{2\pi R}
\sum_{|n|\le N_{\rm coh}}
\exp\left[
-\frac{n^2}{(N_{\rm coh}+1)^2}
\log\frac{1}{\varepsilon_{\rm adm}}
\right]
e^{in(\theta-\theta')}.
\]

The effective mode scale is determined by
\[
\exp\left(-\tau_{\rm adm}\frac{n^2}{R^2}\right)\sim e^{-1}
\quad\Longleftrightarrow\quad
|n|_{\rm eff}\sim R\,\tau_{\rm adm}^{-1/2}
=\frac{R}{\sqrt{\tau_{\rm adm}}}.
\]
Substituting
\[
\tau_{\rm adm}=\frac{R^2}{(N_{\rm coh}+1)^2}\log\frac{1}{\varepsilon_{\rm adm}},
\]
gives the dimensionless retained-mode scale
\[
|n|_{\rm eff}
\sim
\frac{N_{\rm coh}+1}{\sqrt{\log(1/\varepsilon_{\rm adm})}}.
\]

This example shows the selection mechanism explicitly. Once \(R\), \(N_{\rm coh}\), and \(\varepsilon_{\rm adm}\) are fixed, the kernel is fixed. No Gaussian width is chosen by hand.

\section{Conditional Application to Propagator Filters}

The immediate link to the MTT-corrected propagator paper is:
\[
\boxed{
\Delta_{\rm adm}(k)
=
\frac{e^{-\tau_{\rm adm}k^2}}{k^2+m^2},
\qquad
\tau_{\rm adm}
=
\lambda_\ast^{-1}\log\frac{C_Q}{\varepsilon_{\rm adm}}.
}
\]
Thus the exponential factor is inherited from the declared damping rule rather than chosen at
the propagator stage. This remains a conditional Euclidean filter until the same physical
source theorem selects the stabilization generator and its coupling to the propagating field.

For a scalar Euclidean sector with ordinary propagator
\[
\Delta_0(k)=\frac{1}{k^2+m^2},
\]
the admissible MTT-corrected propagator becomes
\[
\boxed{
\Delta_{\rm adm}(k)
=
\frac{e^{-\tau_{\rm adm}k^2}}{k^2+m^2}
}
\]
in the flat-sector model.

The leading low-momentum correction is
\[
\Delta_{\rm adm}(k)
=
\frac{1}{k^2+m^2}
-
\tau_{\rm adm}\frac{k^2}{k^2+m^2}
+
O(\tau_{\rm adm}^2k^4).
\]
Thus an observable probing momentum scale \(E\) is insensitive to the correction when
\[
\tau_{\rm adm}E^2\ll1.
\]
Equivalently,
\[
E\ll \Lambda_{\rm eff}\sim\tau_{\rm adm}^{-1/2}.
\]

\section{Phenomenological Reading}

In natural units \(\hbar=c=1\), \(\tau_{\rm adm}\) has units of inverse mass squared in flat momentum-space sectors:
\[
[\tau_{\rm adm}]=E^{-2}.
\]
The length scale associated with the kernel is
\[
\ell_{\rm coh}\sim \sqrt{\tau_{\rm adm}}.
\]
If an experiment tests locality or propagator behavior up to energy \(E_{\rm max}\) with fractional tolerance \(\eta_{\rm exp}\), then a rough necessary condition is
\[
\tau_{\rm adm}E_{\rm max}^2\lesssim \eta_{\rm exp}.
\]
Hence
\[
\Lambda_{\rm eff}
=
\tau_{\rm adm}^{-1/2}
\gtrsim
\frac{E_{\rm max}}{\sqrt{\eta_{\rm exp}}}.
\]
A real bound requires matching the filtered propagator into the relevant scattering amplitude and observable. The inequality above is only the universal scaling estimate.

\section{What This Resolves and What It Leaves Open}

The criticism of the earlier dictionary papers was correct: showing
\[
K_\epsilon\to\delta
\]
does not identify which \(K_\epsilon\) nature selects.

The present paper supplies a conditional fixed-point execution rule:
\[
(A,P,\chi,C_Q,\lambda_\ast,\varepsilon_{\rm adm})
\Longrightarrow
\tau_{\rm adm}
\Longrightarrow
K_{\rm adm}.
\]
Thus the kernel is unique inside a fully specified regime. The remaining source problem is
not hidden: MTT geometry must select the regime data, including the tolerance, before the
kernel can be called a physical prediction.

\section{Next Execution Targets}

The next tasks are sector-specific:

\begin{enumerate}
\item derive \(A\), \(P\), \(\chi\), \(\lambda_\ast\), and \(\varepsilon_{\rm adm}\) for a concrete MTT sector;
\item compute \(\tau_{\rm adm}\);
\item insert the resulting kernel into propagators, loops, measurement effects, or spectral widths;
\item compare the correction scale to measured data.
\end{enumerate}

The most plausible first physical sector is the neutrino soft sector, because the existing corpus already contains local nil-basin structure, soft-mode Hessians, ordering branches, and measured mass-splitting inputs.

\section{Conclusion}

This paper converts one part of the delta-projection dictionary into a conditional execution
rule. In the fixed-point regime, supplied damping data, spectral gap, admissibility tolerance,
coherent projection, and spectral acceptance determine a unique finite filter.

The central formula is
\[
\boxed{
\tau_{\rm cert}
=
\frac{1}{\lambda_\ast}\log\frac{C_Q}{\varepsilon_{\rm adm}}
}
\]
and the corresponding kernel is
\[
\boxed{
K_{\tau}(x,y)
=
\big\langle x\big|P\chi(A)e^{-\tau A}\chi(A)P\big|y\big\rangle.
}
\]
This gives the first rigorous answer to the reverse problem:
\[
\boxed{
\text{declared fixed-point data}
\Rightarrow
K_{\tau}
\Rightarrow
\text{finite correction}.
}
\]

The missing first-principles step is the selected source map from MTT geometry to those
declared data.

\begin{thebibliography}{9}

\bibitem{EngelNagel}
K.-J. Engel and R. Nagel,
\emph{One-Parameter Semigroups for Linear Evolution Equations},
Graduate Texts in Mathematics 194, Springer (2000).
\href{https://doi.org/10.1007/b97696}{doi:10.1007/b97696}.

\bibitem{Davies}
E. B. Davies,
\emph{Heat Kernels and Spectral Theory},
Cambridge Tracts in Mathematics 92, Cambridge University Press (1989).
\href{https://doi.org/10.1017/CBO9780511566158}{doi:10.1017/CBO9780511566158}.

\end{thebibliography}

\end{document}
